\section{问题重述}

\subsection{问题背景}

无线电干扰源的定位与清除是无线电频谱管理的重要任务。完整的定位流程通常分为初步监测、区域性排查和精准定位三个阶段。前两个阶段依靠固定监测站、移动监测车或无人机等设备进行大范围信号采集与分析，确定干扰源目标的分布区域；但受环境遮挡、低功率干扰源等因素限制，每个干扰源的精准定位最终仍需技术人员在目标分布区域内徒步完成。为尽快消除干扰造成的不良影响，这一阶段耗时越短越好。

某公司拟研发搭载于机器狗的干扰源自动定位清除系统，以便在危险环境中代替技术人员作业。系统依靠便携式无线电方位监测测向机（简称测向机）获取干扰源相对检测点的方位信息，并在足够接近时启用光学探测仪精确定位、以激光枪清除干扰源。

题目给出如下基本条件。

\begin{itemize}[label=\textbullet,leftmargin=2em]
  \item 目标分布区域为半径 1800 米的圆形区域，以圆心为原点、正东为 $x$ 轴正向、正北为 $y$ 轴正向，单位为米；区域中干扰源个数未知。
  \item 每个干扰源的频道互不相同且不随时间变化，频道位于常见 20 个频道之中，依次编号 1 至 20；不同干扰源的信号互不影响。
  \item 干扰源分全向与定向两类。全向源有效覆盖角度范围为 $360^\circ$；定向源具有特定发射方向，有效覆盖角度为定向方向两侧各 $90^\circ$（含），范围之外没有信号辐射。两类源在其有效覆盖范围内信号均匀辐射，场强仅随距离衰减，与角度无关。
  \item 测向机在检测点测得的示向度定义为 $x$ 轴正向逆时针旋转至检测点指向干扰源向量方向的角度，范围为 $[0^\circ,360^\circ)$。受环境与精度限制，读数存在误差；在同一地点、同一段时间内电磁环境干扰固定，重复检测不改变检测误差，只有更换地点后误差才呈现统计规律。从全局看，全部误差落在 $[-1^\circ,1^\circ]$ 内。
  \item 各干扰源的最大可被有效接收距离（有效接收半径）在 1000 米到 1500 米之间。
  \item 测向机不能同时检测多个频道；任意两个频道之间的切换时间为 1 秒。移动过程中无法进行有效检测，必须停下来检测：调整天线朝向并获得稳定读数（或确认无信号）的时间为 5 秒。机器狗沿直线行进，移动速度为 5 米/秒。
  \item 忽略机器狗所携带测向机与干扰源的高度差异，认为二者位于同一平面。
  \item 光学探测仪在距干扰源不超过 20 米时方可精确定位并启动激光枪清除；光学精确定位耗时 3 秒，激光枪清除耗时 2 秒。光学操作是否成功仅与清除点到干扰源的距离有关，与干扰源信号有效覆盖角度范围无关。
  \item 当检测点与干扰源距离不超过 5 米且位于其信号有效覆盖角度范围内时，信号过强而无法获得示向度，此时可直接进行光学精确定位并清除。
\end{itemize}

\subsection{问题重述}

题目要求建立数学模型解决以下四个问题。

\textbf{问题一。} 已知若干检测点坐标及某干扰源关于这些检测点的示向度，给出采用交会定位法计算多边形定位区域直径（即区域内任意两点间距离的最大值）的算法。并回答：以定位区域直径为直径的圆能否覆盖此定位区域？

\textbf{问题二。} 已知某全向干扰源在一个检测点处测得的示向度，给出第二个检测点的选择策略，以获得关于该干扰源较好的定位效果，并给出第二个检测点的候选区域。

\textbf{问题三。} 已知目标区域内干扰源均为全向源，总数在 10 至 16 个之间但具体个数未知，一只机器狗从原点出发、测向机初始频道为 1，需对区域内干扰源自动定位并清除。要求在确保所有干扰源被清除的前提下，使完成任务的时间越短越好。需制定机器狗自动搜索定位及清除的策略与算法，并通过模拟器演练测试，统计"被清除干扰源个数的比例"与"平均定位清除时间"两个指标；其中定位清除总时间包含机器狗移动时间、频道切换时间、检测时间、光学精确定位时间及清除时间。在充分演练后再进行三次正式测试，并按题给表 1 的格式给出每次正式测试的结果（测试案例编码、清除干扰源个数、平均定位清除时间、程序运行时间），同时导出三次正式测试的日志文件放入支撑材料。

\textbf{问题四。} 目标区域中既有全向干扰源又有定向干扰源，总数仍为 10 至 16 个但具体个数未知，定向干扰源的个数及定向方向也未知，其他条件与问题三相同。需在此条件下给出检测与清除的策略及算法，通过模拟器演练测试检验有效性，并同样进行三次正式测试、按表 1 格式给出结果、导出日志文件。

\subsection{各问之间的关系}

四个问题并非彼此独立，而是逐层递进、共享同一套几何基础设施。

问题一提供的是\textbf{定位区域}这一基本对象，以及由它导出的两个量：区域直径 $D(P)$ 与最小包围圆半径 $R(P)$。后两问反复使用的"保证清除"判据正来自 $R(P)\le20$。问题二在问题一的角域表示之上，新增了\textbf{保证接收}这一约束——因为要获得第二次示向度，前提是第二个检测点仍在该源的有效接收半径之内。问题三把前两问的单源、单步工具嵌入多源整场调度，并新增两项前两问没有的内容：负观测的语义（无信号意味着什么）与\textbf{任务终止依据}（凭什么判定已经没有遗漏的干扰源）。问题四则指出问题三所依赖的"全向"前提失效，进而更换无信号解释与覆盖证书，但正观测条件下的定位、包围圆与清除判据仍可原样继承。

本文在写作中严格区分四类陈述：题目给定的事实、由题设推导的结论、为求解而作的建模约定，以及在限定测试范围内观察到的数值结果。第四类仅作为数值支持，不升级为对任意合法场景的保证。
